\section{FORMAL RESTATEMENT}
\textbf{Erd\H{o}s \#4; Lean 4 formalization of the final reduction step.}
Write $\log_j x$ for the $j$-fold iterated natural logarithm and let $p_n$ be the $n$-th prime (zero-indexed in Lean: $p_0=2$). The question asks: is it true that for every $C>0$ there are infinitely many $n$ with
\[
 p_{n+1}-p_n > C\,\frac{\log n\,\log_2 n\,\log_4 n}{(\log_3 n)^2}\ ?
\]
The Lean file below states this as \texttt{Erdos4} and as \texttt{standard\_statement\_of\_stronger\_bound}, in the notation of the formal-conjectures statement.

\section{QUICK LITERATURE/CONTEXT CHECK}
The problem is resolved in the literature (Maynard; Ford--Green--Konyagin--Tao; and jointly Ford, Green, Konyagin, Maynard, Tao, \emph{Long gaps between primes}, J.\ Amer.\ Math.\ Soc.\ 31 (2018), 65--105, \url{https://arxiv.org/abs/1412.5029}). See \url{https://www.erdosproblems.com/4}. No novelty is claimed. The contribution here is a machine-checked formalization of the last implication: the stronger prime-index bound implies the arbitrary-constant statement.

\section{ATTACK PLAN}
Take the stronger bound
\[
 \exists c>0:\quad p_{n+1}-p_n \ge c\,\frac{\log n\,\log_2 n\,\log_4 n}{\log_3 n}\ \text{for infinitely many } n
\]
as an explicit hypothesis (\texttt{StrongerPrimeIndexBound}). The stronger scale equals $\log_3 n$ times the Rankin scale. Since $\log_3 n\to\infty$, a single constant $c$ can be amplified to beat any $C$. Separately, prove the elementary interval lemmas used in the covering-to-gap step of the CRT construction.

\section{WORK}
Full Lean 4 source (\texttt{Erdos4Partial.lean}), unedited:
\begin{verbatim}
import Mathlib.Analysis.SpecialFunctions.Log.Basic
import Mathlib.Data.Nat.Nth
import Mathlib.Data.Nat.Prime.Basic
import Mathlib.Order.Interval.Finset.Basic
import Mathlib.Tactic

/-!
# Erdős problem 4: statement and conditional reduction

STATUS: PARTIAL. This is not an unconditional proof of Erdős problem 4.

Validation: compiled with Lean v4.35.0-rc3 and Mathlib revision
5e0c4e5239cb0a2d86d68a884bf52cfd963fce22.
Run in that Mathlib checkout: lake env lean Erdos4Partial.lean
The printed axiom dependencies are only Lean standard foundations:
propext, Classical.choice, Quot.sound; there are no custom axioms.

Source: the supplied ProblemNumber4(1).pdf, especially equation (6).
References:
* https://www.erdosproblems.com/4
* Ford, Green, Konyagin, Maynard, Tao, Long gaps between primes,
  J. Amer. Math. Soc. 31 (2018), 65–105; https://arxiv.org/abs/1412.5029
* https://github.com/mehmetmars7/Erdosproblems-llm-hunter/blob/main/CONTRIBUTING.md

We prove the last implication in the PDF: its stronger prime-index bound
implies the arbitrary-constant bound in the question. The stronger bound is
an EXPLICIT HYPOTHESIS, not an axiom or an established Lean theorem here.
The covering theorem, CRT construction, and asymptotic passage to equation
(6) are not formalized here. Two elementary interval lemmas are also proved.

Conventions: Nat.nth Nat.Prime is zero-indexed (p 0 = 2), as in the
formal-conjectures statement. The gap uses natural subtraction followed by
coercion to the reals. Real.log is total; eventual positivity is proved,
so the finitely many small indices cause no problem.
-/

noncomputable section
open Filter

namespace Erdos4Partial

def iterLog : ℕ → ℝ → ℝ
  | 0, x => x
  | k + 1, x => Real.log (iterLog k x)

def prime (n : ℕ) : ℕ := Nat.nth Nat.Prime n

def gap (n : ℕ) : ℝ := ((prime (n + 1) - prime n : ℕ) : ℝ)

def rankinScale (n : ℕ) : ℝ :=
  iterLog 1 n * iterLog 2 n * iterLog 4 n / (iterLog 3 n) ^ 2

def strongerScale (n : ℕ) : ℝ :=
  iterLog 1 n * iterLog 2 n * iterLog 4 n / iterLog 3 n

/-- Full target statement; defining a proposition does not prove it. -/
def Erdos4 : Prop :=
  ∀ C : ℝ, 0 < C → {n : ℕ | C * rankinScale n < gap n}.Infinite

/-- The substantive external dependency: equation (6) of the PDF. -/
def StrongerPrimeIndexBound : Prop :=
  ∃ c : ℝ, 0 < c ∧ {n : ℕ | c * strongerScale n ≤ gap n}.Infinite

/-- Every fixed iterated logarithm tends to infinity. -/
theorem iterLog_tendsto (k : ℕ) :
    Tendsto (fun n : ℕ => iterLog k (n : ℝ)) atTop atTop := by
  induction k with
  | zero => simpa [iterLog] using
      (tendsto_natCast_atTop_atTop : Tendsto (fun n : ℕ => (n : ℝ)) atTop atTop)
  | succ k ih => exact Real.tendsto_log_atTop.comp ih

theorem rankinScale_eventually_pos :
    ∀ᶠ n : ℕ in atTop, 0 < rankinScale n := by
  have h1 := (iterLog_tendsto 1).eventually (eventually_gt_atTop (0 : ℝ))
  have h2 := (iterLog_tendsto 2).eventually (eventually_gt_atTop (0 : ℝ))
  have h3 := (iterLog_tendsto 3).eventually (eventually_gt_atTop (0 : ℝ))
  have h4 := (iterLog_tendsto 4).eventually (eventually_gt_atTop (0 : ℝ))
  filter_upwards [h1, h2, h3, h4] with n h1 h2 h3 h4
  exact div_pos (mul_pos (mul_pos h1 h2) h4) (sq_pos_of_pos h3)

/-- Exact algebra behind the extra log-log-log factor. -/
theorem strongerScale_eq (n : ℕ) (h : iterLog 3 (n : ℝ) ≠ 0) :
    strongerScale n = iterLog 3 n * rankinScale n := by
  unfold strongerScale rankinScale
  field_simp [h]

/-- Generic amplification: an unbounded factor turns one constant into any constant. -/
theorem amplify {g s L : ℕ → ℝ} {c : ℝ}
    (hc : 0 < c)
    (hs : ∀ᶠ n in atTop, 0 < s n)
    (hL : Tendsto L atTop atTop)
    (hbig : {n : ℕ | c * (L n * s n) ≤ g n}.Infinite) :
    ∀ C : ℝ, 0 < C → {n : ℕ | C * s n < g n}.Infinite := by
  intro C _hC
  have hthreshold : ∀ᶠ n in atTop, C / c < L n :=
    hL.eventually (eventually_gt_atTop (C / c))
  obtain ⟨N, hN⟩ := eventually_atTop.1 (hs.and hthreshold)
  apply Set.infinite_iff_exists_gt.2
  intro M
  obtain ⟨n, hn, hlarge⟩ := hbig.exists_gt (max M N)
  have hMn : M < n := lt_of_le_of_lt (le_max_left M N) hlarge
  have hNn : N ≤ n := le_of_lt (lt_of_le_of_lt (le_max_right M N) hlarge)
  obtain ⟨hsn, hLn⟩ := hN n hNn
  have hCL : C < c * L n := by
    have := (div_lt_iff₀ hc).1 hLn
    nlinarith
  refine ⟨n, ?_, hMn⟩
  change C * s n < g n
  calc
    C * s n < (c * L n) * s n := mul_lt_mul_of_pos_right hCL hsn
    _ = c * (L n * s n) := mul_assoc _ _ _
    _ ≤ g n := hn

/-- CONDITIONAL result. The stronger prime-index theorem remains a hypothesis. -/
theorem erdos4_of_stronger_bound (hstrong : StrongerPrimeIndexBound) : Erdos4 := by
  obtain ⟨c, hc, hbig⟩ := hstrong
  have hidentity (n : ℕ) : strongerScale n = iterLog 3 n * rankinScale n := by
    by_cases hn : iterLog 3 (n : ℝ) = 0
    · simp [strongerScale, rankinScale, hn]
    · exact strongerScale_eq n hn
  have hbig' : {n : ℕ | c * (iterLog 3 n * rankinScale n) ≤ gap n}.Infinite := by
    simpa only [hidentity] using hbig
  exact amplify hc rankinScale_eventually_pos (iterLog_tendsto 3) hbig'

/-- The normalization agrees with the displayed formula, including multiplication order. -/
theorem rankin_formula (C : ℝ) (n : ℕ) :
    C * rankinScale n =
      C * Real.log (Real.log (n : ℝ)) *
        Real.log (Real.log (Real.log (Real.log (n : ℝ)))) /
        (Real.log (Real.log (Real.log (n : ℝ)))) ^ 2 * Real.log (n : ℝ) := by
  simp only [rankinScale, iterLog]
  ring

/-- The conditional theorem in the notation used by formal-conjectures. -/
theorem standard_statement_of_stronger_bound (hstrong : StrongerPrimeIndexBound) :
    ∀ C : ℝ, 0 < C →
      {n : ℕ | (↑(Nat.nth Nat.Prime (n + 1) - Nat.nth Nat.Prime n) : ℝ) >
        C * Real.log (Real.log (n : ℝ)) *
          Real.log (Real.log (Real.log (Real.log (n : ℝ)))) /
          (Real.log (Real.log (Real.log (n : ℝ)))) ^ 2 * Real.log (n : ℝ)}.Infinite := by
  intro C hC
  have h := erdos4_of_stronger_bound hstrong C hC
  simpa only [rankin_formula, gap, prime] using h

/-- A proper divisor at least two rules out primality, including q = a+t. -/
theorem translated_not_prime {a t q : ℕ}
    (hq : 2 ≤ q) (hqa : q ≤ a) (ht : 0 < t) (hdiv : q ∣ a + t) :
    ¬ Nat.Prime (a + t) := by
  intro hp
  exact (Nat.prime_def_lt'.1 hp).2 q hq (by omega) hdiv

/-- Once translated divisibility is established, a covered interval is prime-free. -/
theorem covered_interval_prime_free {a y : ℕ}
    (hcover : ∀ t : ℕ, 1 ≤ t → t ≤ y →
      ∃ q : ℕ, 2 ≤ q ∧ q ≤ a ∧ q ∣ a + t) :
    ∀ t : ℕ, 1 ≤ t → t ≤ y → ¬ Nat.Prime (a + t) := by
  intro t ht hty
  obtain ⟨q, hq, hqa, hdiv⟩ := hcover t ht hty
  exact translated_not_prime hq hqa (by omega) hdiv

/-- A prime to the right of a prime-free interval gives the required gap. -/
theorem gap_of_prime_free_interval {a y p p' : ℕ}
    (hp : p ≤ a) (hp' : a < p') (hprime : Nat.Prime p')
    (hfree : ∀ t : ℕ, 1 ≤ t → t ≤ y → ¬ Nat.Prime (a + t)) :
    y < p' - p := by
  have hout : a + y < p' := by
    by_contra h
    have ht : 1 ≤ p' - a := by omega
    have hty : p' - a ≤ y := by omega
    have heq : a + (p' - a) = p' := by omega
    exact (hfree (p' - a) ht hty) (heq.symm ▸ hprime)
  omega

#print axioms erdos4_of_stronger_bound
#print axioms covered_interval_prime_free
#print axioms gap_of_prime_free_interval

end Erdos4Partial
\end{verbatim}

Main results:
\begin{itemize}
\item \texttt{iterLog\_tendsto}: every fixed iterated logarithm tends to infinity.
\item \texttt{amplify}: if $c\,L(n)s(n)\le g(n)$ infinitely often, $s>0$ eventually and $L\to\infty$, then $C\,s(n)<g(n)$ infinitely often for every $C>0$.
\item \texttt{erdos4\_of\_stronger\_bound}: \texttt{StrongerPrimeIndexBound} $\Rightarrow$ \texttt{Erdos4}.
\item \texttt{standard\_statement\_of\_stronger\_bound}: the same conclusion in the displayed formula, with \texttt{rankin\_formula} checking the normalization.
\item \texttt{translated\_not\_prime}, \texttt{covered\_interval\_prime\_free}, \texttt{gap\_of\_prime\_free\_interval}: a covered interval $(a,a+y]$ has no primes, which gives a gap $>y$.
\end{itemize}

\section{VERIFICATION}
Compiled with Lean v4.35.0-rc3 and Mathlib revision \texttt{5e0c4e5239cb0a2d86d68a884bf52cfd963fce22} via \texttt{lake env lean Erdos4Partial.lean}. \texttt{\#print axioms} reports only \texttt{propext}, \texttt{Classical.choice}, \texttt{Quot.sound}. There is no \texttt{sorry} and there are no custom axioms. \texttt{StrongerPrimeIndexBound} appears only as a hypothesis of the theorem, never as an axiom.

\section{FINAL STATUS}
\textbf{PARTIAL.} Proved in Lean: the implication from the stronger prime-index bound to the problem's statement, plus the elementary prime-free-interval lemmas. Not formalized: the Ford--Green--Konyagin--Maynard--Tao covering theorem, the CRT construction, and the asymptotic passage from the covering bound to the stronger prime-index bound (equation (6) of the source notes). The first remaining gap is to prove \texttt{StrongerPrimeIndexBound} in Lean. Subject to expert review.
